\section{Related Work}
\label{sec:related}

\subsection{Large Language Models for Automated Test Suite Generation}
Recent advancements in Large Language Models (LLMs) have transformed automated software testing from heuristic code generation to semantic test synthesis. Early generative approaches utilized pre-trained code models such as Codex and StarCoder to generate unit tests from method signatures and docstrings~\cite{bareiss2022code}. Subsequent frameworks introduced few-shot prompting and context-augmented generation to improve assertion accuracy and line coverage. For instance, ChatUnitTest~\cite{xie2023chatunitest} leveraged ChatGPT to synthesize Java unit tests, achieving up to 74.2\% line coverage across open-source benchmarks. Similarly, CODAMOSA~\cite{lemieux2023codamosa} combined LLMs with search-based testing to escape coverage plateaus, increasing branch coverage by 18.5\% compared to traditional generators. In repository-level contexts, DocPrompting~\cite{docprompting} and RepoBench~\cite{repobench} incorporated documentation and multi-file context to guide code completion.

Despite these achievements, a shared limitation of static LLM-based test generators is their vulnerability to \textit{compliance bias}~\cite{testexplora}. Because static models generate test code in a single inference pass without execution signals, they exhibit a passive tendency to conform to existing code implementations. When evaluated on defective repositories, static generators generate unit tests that pass on buggy code, yielding a Fail-to-Pass (F2P) rate floor of only 16.06\% on the TestExplora benchmark~\cite{testexplora}.

\subsection{Search-Based Software Testing and Execution-Driven Agents}
Search-Based Software Testing (SBST) tools such as EvoSuite~\cite{evosuite} for Java and Randoop~\cite{randoop} for object-oriented systems utilize genetic algorithms and feedback-directed random input generation to maximize structural coverage. SBST tools excel at finding edge-case crashes and executing code paths dynamically; however, they lack natural language understanding, often producing unreadable test inputs with weak or missing semantic assertions~\cite{siddiq2023exploratory, schafer2023empirical}.

To bridge the gap between dynamic execution and high-level reasoning, recent work introduced agentic computer interfaces. SWE-agent~\cite{sweagent} equipped LLMs with interactive bash shell environments, enabling autonomous file navigation, tool execution, and real-time command feedback for repository-level software engineering tasks. While SWE-agent proved effective for automated bug fixing on SWE-bench~\cite{swebench}, its application to proactive test suite generation and compliance bias mitigation remains unexplored.

\subsection{Positioning of This Work}
Unlike prior work, this paper is the first to evaluate interactive terminal execution feedback using a multi-agent framework (combining DeepSeek-V3 and Llama-3.3-70B via the SWE-agent CLI) for proactive repository-level test suite generation on the TestExplora benchmark to break compliance bias.
